\documentclass[11pt]{article}
% Shared preamble for per-Python-file documentation (input via \input{...}).
\usepackage[margin=1in]{geometry}
\usepackage[T1]{fontenc}
\usepackage[utf8]{inputenc}
\usepackage{lmodern}
\usepackage{hyperref}
\usepackage{xcolor}
\usepackage{listings}
\usepackage{listingsutf8}
\lstset{
  basicstyle=\ttfamily\small,
  columns=fullflexible,
  breaklines=true,
  upquote=true,
  inputencoding=utf8,
  extendedchars=true,
  frame=single,
  framerule=0.2pt,
  tabsize=2
}


\title{\texttt{model/\_\_init\_\_.py} (Q\&A)}
\author{}
\date{}

\begin{document}
\maketitle

\paragraph{Q: What is the purpose of \texttt{model/\_\_init\_\_.py}?}
\textbf{A:} It turns the \lstinline|model/| folder into an importable Python package and provides a single import surface that re-exports the main model classes used by the training scripts.

\paragraph{Q: Which models are re-exported?}
\textbf{A:} It re-exports \lstinline|WideAndDeep|, \lstinline|Cread|, \lstinline|D2Q|, \lstinline|TPM|, and \lstinline|EGMN| so scripts can write \lstinline|from model import EGMN| instead of importing individual modules.

\paragraph{Q: Why does this pattern matter for the repo?}
\textbf{A:} The run scripts under the repo root depend on stable, short imports; re-exporting classes here decouples script imports from the file layout (e.g. renaming \lstinline|model/egmn.py| would only require updating this file and not every run script).

\paragraph{Q: Full source code?}
\textbf{A:}
\lstinputlisting[language=Python]{/workspace/model/__init__.py}

\end{document}
